\documentclass[11pt]{article}
\usepackage[margin=1in]{geometry}
\usepackage{amsmath,amssymb,amsthm}
\usepackage[hidelinks]{hyperref}
\newtheorem{theorem}{Theorem}
\newtheorem{lemma}[theorem]{Lemma}
\newtheorem{corollary}[theorem]{Corollary}
\theoremstyle{remark}
\newtheorem{remark}[theorem]{Remark}
\DeclareMathOperator{\Fail}{Fail}
\DeclareMathOperator{\rad}{rad}
\DeclareMathOperator{\lcm}{lcm}
\title{Prime support and growth of Stirling repair factors}
\author{}
\date{September 8, 2026}
\begin{document}
\maketitle
\begin{abstract}
For fixed $k\geq1$, let $F_k$ be the least positive integer such that
$(F_k S(n+k-1,k))_{n\geq1}$ counts the periodic points of a map, where $S(N,k)$
is a Stirling number of the second kind. We prove
\[
 \rad((k-1)!)\mid F_k\mid\lcm(1,\ldots,k-1).
\]
Thus the prime divisors of $F_k$ are exactly the primes below $k$, and the
previous factorial bound can be replaced by a least common multiple.
We also prove $v_p(F_k)=1$ whenever $p<k\leq p^2$, and
$v_p(F_{p^j+1})=j$ for every prime $p$ and $j\geq1$.
A formula for prime-power values in terms of the base-$p$ digits of $k-1$
also yields
$\log(F_k/\rad((k-1)!))\sim\sqrt{k}$.
These results establish Conjectures 1, 2, 3, and 4(d) of Miska and Ward.
The upper bound follows from a Frobenius congruence for an integral recurrence;
the prime-support obstruction comes from the repeated poles of its generating
function modulo $p$.
\end{abstract}

\section{Introduction}
An integer sequence $(a_n)_{n\geq1}$ satisfies the \emph{Dold condition} if
\begin{equation}\label{eq:dold}
 n\mid\sum_{d\mid n}\mu(n/d)a_d\qquad(n\geq1).
\end{equation}
For a nonnegative integer sequence, realizability as the numbers of fixed
points of the iterates of a map is equivalent to \eqref{eq:dold} together with
nonnegativity of the sums in \eqref{eq:dold}. Indeed, their quotients by $n$
are precisely the required numbers of orbits of length $n$. A disjoint union
of those finite cycles supplies a realizing map.

Miska and Ward introduced the least positive multiplier $\Fail(a)$ making a
sequence realizable and studied it for Stirling sequences
\cite{MW}. Throughout this paper,
\[
 a_n^{(k)}=S(n+k-1,k),\qquad
 F_k=\Fail\bigl((a_n^{(k)})_{n\geq1}\bigr).
\]
They proved that $F_k$ is finite and divides $(k-1)!$
\cite[Theorem~2.2]{MW}. Their Conjectures~1 and~2 predict, respectively, that
every prime below $k$ divides $F_k$, and that primes in
$[\sqrt{k},k)$ have valuation exactly one. Their Conjecture~4(d) predicts
$v_p(F_{p^j+1})=j$. Their Conjecture~3 predicts that
$F_k/\rad((k-1)!)$ tends to infinity. We prove all four assertions, improve
the factorial upper bound, and determine the logarithmic growth of that ratio.

\begin{theorem}\label{thm:main}
For every integer $k\geq1$,
\begin{equation}\label{eq:bounds}
 \rad((k-1)!)\mid F_k\mid L_{k-1},
 \qquad L_{k-1}=\lcm(1,\ldots,k-1),
\end{equation}
where $L_0=1$. More precisely, for every prime $p<k$,
\[
 1\leq v_p(F_k)\leq\lfloor\log_p(k-1)\rfloor.
\]
In addition:
\begin{enumerate}
\item if $p<k\leq p^2$, then $v_p(F_k)=1$;
\item if $k=p^j+1$ with $j\geq1$, then $v_p(F_k)=j$.
\end{enumerate}
\end{theorem}

\begin{theorem}\label{thm:growth}
Let $\vartheta(x)=\sum_{p\leq x}\log p$, where the sum is over primes.
As $k\to\infty$, with $K=k-1$,
\begin{equation}\label{eq:theta-growth}
 \log\frac{F_k}{\rad(K!)}
 =\vartheta(\sqrt K)+O\bigl(K^{1/3}(\log K)^2\bigr).
\end{equation}
Consequently,
\begin{equation}\label{eq:growth}
 \log\frac{F_k}{\rad((k-1)!)}\sim\sqrt{k}.
\end{equation}
\end{theorem}

The estimate \eqref{eq:theta-growth} is elementary. Only the passage to
\eqref{eq:growth} uses the prime number theorem.

The upper bound is sharp at the prime $p$ whenever $k=p^j+1$.
The argument for prime support applies to the entire sequence and does not
require locating an individual prime-power witness. In particular, our result
does not establish the proposed finite cutoff for computing $F_k$ described
in \cite[Section~4.5]{BGW}.

\section{Local congruences and the integral recurrence}
For a sequence $a$ write $D_n(a)=\sum_{d\mid n}\mu(n/d)a_d$.
The following standard local form of the Dold condition will be useful; we
include the proof to fix the indexing.

\begin{lemma}\label{lem:local}
An integer sequence $a$ satisfies the Dold condition if and only if
\begin{equation}\label{eq:local}
 a_{mp^r}\equiv a_{mp^{r-1}}\pmod {p^r}
\end{equation}
for every prime $p$, every $r\geq1$, and every $m\geq1$ with $p\nmid m$.
\end{lemma}
\begin{proof}
M\"obius inversion gives $a_n=\sum_{d\mid n}D_d(a)$. Thus
\[
 a_{mp^r}-a_{mp^{r-1}}=\sum_{d\mid m}D_{dp^r}(a),
\]
which proves one implication. Conversely,
\[
 D_{mp^r}(a)=\sum_{d\mid m}\mu(m/d)
       \bigl(a_{dp^r}-a_{dp^{r-1}}\bigr).
\]
The congruences imply $p^r\mid D_{mp^r}(a)$. Applying this at each prime
factor of an arbitrary index gives \eqref{eq:dold}.
\end{proof}

\begin{lemma}\label{lem:sign}
For every $k\geq1$, the sequence $a^{(k)}$ satisfies $D_n(a^{(k)})\geq0$
for all $n\geq1$. Consequently, $F_k$ is its least positive Dold multiplier.
\end{lemma}
\begin{proof}
For $k=1$ the sequence is constant equal to one, so the assertion is
immediate. For $k\geq2$, the Stirling recurrence gives
\[
 a_{n+1}^{(k)}=k a_n^{(k)}+S(n+k-1,k-1)\geq2a_n^{(k)}.
\]
As $a_1^{(k)}=1$, for $n\geq2$ we have
\[
 \sum_{d=1}^{n-1}a_d^{(k)}
 \leq a_n^{(k)}\sum_{d=1}^{n-1}2^{d-n}<a_n^{(k)}.
\]
Therefore $D_n(a^{(k)})\geq a_n^{(k)}-\sum_{d<n}a_d^{(k)}>0$.
Positive multiplication preserves this sign condition.
\end{proof}

Fix $k\geq2$, and suppress the superscript on $a_n$. Define
\[
 P_k(z)=\prod_{j=1}^k(1-jz),\qquad
 G_k(X)=\prod_{j=1}^k(X-j).
\]
The usual generating identity for Stirling numbers is
\begin{equation}\label{eq:gf}
 H_k(z):=\sum_{n\geq1}a_nz^{n-1}=\frac1{P_k(z)}.
\end{equation}
It follows directly, for example, by applying the Stirling recurrence to
$\sum_{t\geq0}S(k+t,k)z^t$: this series is the corresponding series for
$k-1$, divided by $1-kz$, and the initial series for $k=1$ is $(1-z)^{-1}$.

\begin{lemma}\label{lem:companion}
Set $a_0=S(k-1,k)=0$. There are an integer $k\times k$ matrix $T$,
a vector $v\in\mathbb Z^k$, and an integer linear functional $\lambda$
such that
\[
 G_k(T)=0,\qquad a_n=\lambda(T^n v)\quad(n\geq0).
\]
\end{lemma}
\begin{proof}
Write $G_k(X)=X^k+c_1X^{k-1}+\cdots+c_k$.
Multiplying \eqref{eq:gf} by $P_k(z)=1+c_1z+\cdots+c_kz^k$ shows
\[
 a_{n+k}+c_1a_{n+k-1}+\cdots+c_ka_n=0\qquad(n\geq0).
\]
For $n=0$ this uses the coefficient of $z^{k-1}$, which is zero since
$k\geq2$, and the fact that $a_0=0$. For $n\geq1$ it uses the coefficient
of $z^{n+k-1}$. Take the companion transition matrix for the states
$(a_n,a_{n+1},\ldots,a_{n+k-1})^{\mathsf T}$, take the initial state as $v$,
and let $\lambda$ extract its first coordinate. The companion identity
$G_k(T)=0$ gives the result.
\end{proof}

The integral initialization at zero is useful: the matrix powers represent
$a_n$ itself, without a shift in the exponent.

\section{The least-common-multiple upper bound}
\begin{lemma}\label{lem:lift}
Let $X,Y$ be commuting integer matrices. If $X\equiv Y\pmod {p^a}$,
where $p$ is prime and $a\geq1$, then
\[
 X^{p^t}\equiv Y^{p^t}\pmod {p^{a+t}}\qquad(t\geq0).
\]
\end{lemma}
\begin{proof}
Write $X=Y+p^aC$; then $C$ commutes with $Y$. In the binomial expansion
of $X^p-Y^p$, the terms of degree $i$ in $C$, for $1\leq i<p$, are
divisible by $p^{ai+1}$, hence by $p^{a+1}$. The last term is divisible
by $p^{ap}$, and $ap\geq a+1$. Iteration proves the assertion, including
$p=2$.
\end{proof}

\begin{lemma}\label{lem:frob}
Fix a prime $p$ and put
\[
 m=\lceil k/p\rceil,\qquad h=\lceil\log_p m\rceil.
\]
For the matrix $T$ of Lemma~\ref{lem:companion},
\[
 T^{p^{h+1}}\equiv T^{p^h}\pmod p.
\]
\end{lemma}
\begin{proof}
Every residue class modulo $p$ occurs among $1,\ldots,k$ at most $m$
times. Over $\mathbb F_p$ this implies
$G_k(X)\mid(X^p-X)^m$, so $(T^p-T)^m=0$ modulo $p$.
Since $p^h\geq m$, raising $T^p-T$ to the power $p^h$ gives zero.
The characteristic-$p$ binomial identity for commuting matrices gives
$T^{p^{h+1}}-T^{p^h}=0$ modulo $p$.
\end{proof}

\begin{proof}[Proof of the upper bound in Theorem~\ref{thm:main}]
The case $k=1$ is immediate. For $k\geq2$, fix a prime $p$ and use the
notation of Lemma~\ref{lem:frob}. For $r\geq h+1$, Lemma~\ref{lem:lift}
gives
\[
 T^{p^r}\equiv T^{p^{r-1}}\pmod {p^{r-h}}.
\]
Taking arbitrary positive powers preserves the congruence. Applying
$\lambda(\,\cdot\,v)$ therefore gives
\[
 a_{tp^r}\equiv a_{tp^{r-1}}\pmod {p^{r-h}}
 \qquad(t\geq1,\ r\geq h+1).
\]
Multiplication by $p^h$ supplies the required local modulus $p^r$.
For $1\leq r\leq h$ the same conclusion is automatic after multiplication
by $p^h$.

If $p\geq k$, then $m=1$ and $h=0$. If $p<k$, then
\begin{equation}\label{eq:h}
 h=\lfloor\log_p(k-1)\rfloor.
\end{equation}
Indeed, for the integer on the right, $p^h<k\leq p^{h+1}$, and hence
$p^{h-1}<\lceil k/p\rceil\leq p^h$.
The product of these local multipliers is
\[
 \prod_{p<k}p^{\lfloor\log_p(k-1)\rfloor}=L_{k-1}.
\]
Lemmas~\ref{lem:local} and~\ref{lem:sign} show that this is a realizability
multiplier. The least multiplier divides every multiplier: it is the least
common multiple of the denominators of the rational numbers $D_n(a)/n$
in lowest terms. Consequently $F_k\mid L_{k-1}$.
\end{proof}

\section{Every prime below \texorpdfstring{$k$}{k} is necessary}
\begin{proof}[Proof of the lower bound in Theorem~\ref{thm:main}]
Let $p<k$, and suppose that $p\nmid F_k$. The local congruences for $F_ka$,
reduced modulo $p$, imply
\begin{equation}\label{eq:cartier}
 a_{pn}=a_n\quad\text{in }\mathbb F_p\qquad(n\geq1).
\end{equation}
Here one writes $n=tp^r$ with $p\nmid t$ and applies the local congruence
at level $r+1$, cancelling the nonzero scalar $F_k$ modulo $p$.
Iteration gives $a_{p^h n}=a_n$ for every $h\geq0$.

Choose $h$ as in Lemma~\ref{lem:frob} and set $B=T^{p^h}$ modulo $p$.
Then $B^p=B$, whence
$B^{n+p-1}=B^n$ for every $n\geq1$.
The matrix representation and \eqref{eq:cartier} imply
\[
 a_{n+p-1}=a_n\quad\text{in }\mathbb F_p\qquad(n\geq1).
\]
Thus the generating series $H_k(z)$ modulo $p$ has period $p-1$, and
$(1-z^{p-1})H_k(z)$ is a polynomial. By \eqref{eq:gf}, this forces
\[
 P_k(z)\mid1-z^{p-1}\quad\text{in }\mathbb F_p[z].
\]
There is no numerator cancellation because the numerator in \eqref{eq:gf}
is the constant one. But $k>p$ means that the factors with $j=1$ and
$j=p+1$ both occur, so $(1-z)^2\mid P_k(z)$ modulo $p$.
The polynomial $1-z^{p-1}$ is squarefree: its derivative in characteristic
$p$ is $z^{p-2}$, which has no common root with it. This is a contradiction,
also when $p=2$. Hence $p\mid F_k$ for every $p<k$.
\end{proof}

\section{Prime-power values and exact valuations}
The next formula extends the first-level congruence
$S(p+k-1,k)\equiv\lceil k/p\rceil\pmod p$ of Miska and Ward
\cite[proof of Lemma~2.5]{MW}.

\begin{lemma}\label{lem:digits}
Let $p$ be prime, $k\geq1$, and $r\geq0$. Write
\[
 k-1=\sum_{i\geq0}d_i p^i,\qquad 0\leq d_i<p.
\]
Then
\begin{equation}\label{eq:digits}
 S(p^r+k-1,k)\equiv\prod_{i=1}^r(d_i+1)\pmod p,
\end{equation}
where the empty product is one.
\end{lemma}
\begin{proof}
Write $k=qp+c$, with $0\leq c<p$, and put $R=(p^r-1)/(p-1)$.
Over $\mathbb F_p$ we have
\[
 P_k(z)=(1-z^{p-1})^qP_c(z),\qquad P_0(z)=1.
\]
If $1\leq c<p$, every coefficient of $1/P_c(z)$ at a nonnegative
degree $e$ divisible by $p-1$ equals one modulo $p$. Indeed, the
coefficient is $S(e+c,c)$, and inclusion--exclusion gives
\[
 S(e+c,c)=\frac1{c!}\sum_{j=1}^c(-1)^{c-j}\binom cj j^{e+c}
 \equiv S(c,c)=1\pmod p.
\]
Here $c!$ is invertible modulo $p$ and Fermat's congruence applies to each
$j\in\{1,\ldots,c\}$.
Thus, when $c>0$ and $q\geq1$, the coefficient of $z^{(p-1)R}$ in
$1/P_k(z)$ is
\[
 \sum_{t=0}^R\binom{q+t-1}{t}=\binom{q+R}{R}\pmod p.
\]
When $c>0$ and $q=0$, this coefficient is one, which agrees with the final
binomial expression. If $c=0$, only the degree-zero coefficient of
$1/P_0(z)$ contributes, and the answer is $\binom{q+R-1}{R}$.
Putting
$J=\lceil k/p\rceil-1=\lfloor(k-1)/p\rfloor$ combines both cases into
\begin{equation}\label{eq:binomial}
 a_{p^r}^{(k)}\equiv\binom{J+R}{R}\pmod p.
\end{equation}

The base-$p$ digits of $J$ are $d_1,d_2,\ldots$, and those of $R$ are
one in positions $0,\ldots,r-1$ and zero thereafter.
If all $d_1,\ldots,d_r$ are at most $p-2$, adding $J$ and $R$ produces no
carry. Lucas's coefficient formula therefore evaluates
\eqref{eq:binomial} as $\prod_{i=1}^r(d_i+1)$.
If one of these digits is $p-1$, take the first such digit. There was no
earlier carry, and the corresponding digit of $J+R$ is zero. Its Lucas
factor is $\binom01=0$, and the proposed product is zero as well.
For completeness, the coefficient formula used here follows by comparing
coefficients in
\[
 (1+x)^N=\prod_{i\geq0}(1+x^{p^i})^{N_i}\quad\text{in }\mathbb F_p[x],
\]
where $N_i$ are the base-$p$ digits of $N$.
\end{proof}

\begin{corollary}\label{cor:sharp-digits}
Let $p<k$, put $h=\lfloor\log_p(k-1)\rfloor$, and use the digits in
Lemma~\ref{lem:digits}. If $1\leq r\leq h$, $d_r\ne0$, and
$d_i\ne p-1$ for every $1\leq i<r$, then $v_p(F_k)\geq r$.
In particular, if $d_i\ne p-1$ for every $1\leq i<h$, then
\[
 v_p(F_k)=h.
\]
\end{corollary}
\begin{proof}
Subtracting successive formulas in \eqref{eq:digits} gives
\[
 a_{p^r}^{(k)}-a_{p^{r-1}}^{(k)}
 \equiv d_r\prod_{i=1}^{r-1}(d_i+1)\pmod p.
\]
Under the stated conditions this is nonzero. The local Dold congruence at
level $r$ forces $p^r\mid F_k$. At $r=h$, the highest digit $d_h$ is
nonzero, and the upper bound in Theorem~\ref{thm:main} gives equality.
\end{proof}

\begin{remark}
Lemma~\ref{lem:digits} also supplies an explicit witness for full prime
support. For $p<k$, let $s\geq1$ be the first position with $d_s\ne0$.
Then $a_{p^s}^{(k)}-a_{p^{s-1}}^{(k)}\not\equiv0\pmod p$, with $p^s<k$.
This proves the prime-support assertion by a second method and locates
one necessary prime-power factor below $k$. It does not locate the
maximum possible valuation of the repair factor.
\end{remark}

\begin{proof}[Proof of Theorem~\ref{thm:main}(1) and (2)]
When $p<k\leq p^2$, we have $h=1$, so the digit restrictions in
Corollary~\ref{cor:sharp-digits} are empty and $v_p(F_k)=1$.
When $k=p^j+1$, the digits of $k-1=p^j$ are zero below position $j$
and one at position $j$. The same corollary gives $v_p(F_k)=j$.
In particular, for $p=2$ the corollary also gives
$v_2(F_k)=h$ at both $k=2^h+1$ and $k=2^h+2$, for every $h\geq1$.
\end{proof}

\section{Growth beyond the squarefree factor}
\begin{proof}[Proof of Theorem~\ref{thm:growth}]
Put $K=k-1$ and $E_k=\log(F_k/\rad(K!))$. By Theorem~\ref{thm:main},
\begin{align*}
 0\leq E_k
 &\leq\sum_{j=2}^{\lfloor\log_2 K\rfloor}\vartheta(K^{1/j})\\
 &=\vartheta(\sqrt K)+O\bigl(K^{1/3}(\log K)^2\bigr).
\end{align*}
For the error estimate, there are $O(\log K)$ terms with $j\geq3$, and
each is at most $K^{1/3}\log K$, by the elementary inequality
$\vartheta(x)\leq x\log x$.

For the lower bound, consider primes
\[
 K^{1/3}<p\leq\sqrt K.
\]
For each such prime, $h=\lfloor\log_p K\rfloor=2$.
Corollary~\ref{cor:sharp-digits} gives $v_p(F_k)=2$ unless the digit
$d_1$ in $K=d_2p^2+d_1p+d_0$ is $p-1$.
In that exceptional case, $t=d_2+1$ satisfies
\begin{equation}\label{eq:exception}
 tp^2-p\leq K<tp^2,\qquad 2\leq t\leq K^{1/3}+1.
\end{equation}
For a fixed $t\geq2$, at most one positive integer $p$ can satisfy
\eqref{eq:exception}. Indeed, the lower endpoint for $p+1$ exceeds the
upper endpoint for $p$, since
\[
 t(p+1)^2-(p+1)-tp^2=2tp+t-p-1>0.
\]
The intervals increase with $p$, so they are pairwise disjoint.
Consequently, at most $\lfloor K^{1/3}\rfloor+1$ primes in the specified
range are exceptional.

Every nonexceptional prime contributes at least $\log p$ to $E_k$.
It follows that
\begin{align*}
 E_k&\geq\vartheta(\sqrt K)-\vartheta(K^{1/3})
       -(\lfloor K^{1/3}\rfloor+1)\log\sqrt K\\
    &=\vartheta(\sqrt K)-O(K^{1/3}\log K).
\end{align*}
This proves \eqref{eq:theta-growth}. The classical prime number theorem
$\vartheta(x)\sim x$ \cite{Hadamard}, together with
$K^{1/3}(\log K)^2=o(\sqrt K)$, proves \eqref{eq:growth}.
In particular, $F_k/\rad((k-1)!)\to\infty$.
\end{proof}

\section{Remaining questions}
The exact value of $F_k$ for general $k$ is still not determined by these
arguments. In particular, the finite-cutoff assertion discussed in
\cite[Section~4.5]{BGW} would compute $F_k$ from the denominators of
$D_n(a^{(k)})/n$ only for indices $n$ at most the largest prime power
strictly below $k$. Our upper bound controls all prime-power valuations,
but does not show that their maxima occur within that finite range.

The results above leave open parts (a)--(c) of Conjecture~4 in \cite{MW}. In particular, the assertion
$v_p(F_{p^j})=1$ for $j>1$ requires a sharper upper bound than the one
proved here. Distinguishing those special cancellations from the general
Frobenius bound appears necessary for an exact formula.

\begin{thebibliography}{9}
\bibitem{BGW}
J. Byszewski, G. Graff, and T. Ward,
\emph{Dold sequences, periodic points, and dynamics},
Bull. London Math. Soc. \textbf{53} (2021), 1263--1298.
\href{https://doi.org/10.1112/blms.12531}{doi:10.1112/blms.12531}.

\bibitem{Hadamard}
J. Hadamard,
\emph{Sur la distribution des z\'eros de la fonction $\zeta(s)$ et ses
cons\'equences arithm\'etiques},
Bull. Soc. Math. France \textbf{24} (1896), 199--220.
\href{https://doi.org/10.24033/bsmf.545}{doi:10.24033/bsmf.545}.

\bibitem{MW}
P. Miska and T. Ward,
\emph{Stirling numbers and periodic points},
Acta Arith. \textbf{201} (2021), 421--435.
\href{https://doi.org/10.4064/aa210309-9-8}{doi:10.4064/aa210309-9-8}.
Author preprint: \href{https://arxiv.org/abs/2102.07561}{arXiv:2102.07561}.
\end{thebibliography}
\end{document}
